% Figure: Thermal circuit equivalent analogy
% Chapter: ch19 - Thermal-Electro-Optic Coupling
% Description: Analogy between electrical circuit and thermal circuit for intuitive understanding of thermal resistance

\resizebox{\textwidth}{!}{%
\begin{tikzpicture}[
    scale=0.9,
    elec_style/.style={thick, blue!70!black},
    therm_style/.style={thick, red!70!black},
    component/.style={fill=white, draw, minimum width=1.2cm, minimum height=0.8cm},
    arrow_flow/.style={->, >=stealth, thick}
  ]
    
    % ========== Left: Electrical Circuit (Familiar) ==========
    \begin{scope}[shift={(0,0)}]
      \node[anchor=south, font=\large\bfseries] at (2.5, 4.2) {(a) 电路类比};
      
      % Voltage source
      \draw[elec_style] (0.5, 0.5) to[battery1, l=$V_{\mathrm{dd}}$] (0.5, 3.5);
      
      % Resistor
      \draw[elec_style] (0.5, 3.5) to[R=$R$, l_={$R=\rho L/A$}] (4.5, 3.5);
      
      % Ground
      \draw[elec_style] (4.5, 3.5) -- (4.5, 0.5) node[ground]{};
      
      % Current arrow
      \draw[elec_style, ->] (1.5, 3.7) -- (2.5, 3.7) node[midway, above] {$I$};
      
      % Ohm's law box
      \draw[rounded corners, fill=blue!15] (5.2, 2.5) rectangle (8.5, 3.8);
      \node[anchor=west, font=\small] at (5.4, 3.4) {欧姆定律：};
      \node[anchor=west, font=\normalsize] at (5.4, 2.9) {$V = I \cdot R$};
      
      % Labels
      \node[anchor=east, font=\small] at (0.3, 3.5) {电压源};
      \node[anchor=west, font=\small] at (4.7, 3.5) {负载};
      \node[anchor=north, font=\small] at (2.5, 0.3) {地};
    \end{scope}
    
    % ========== Right: Thermal Circuit (Analogous) ==========
    \begin{scope}[shift={(10,0)}]
      \node[anchor=south, font=\large\bfseries] at (2.5, 4.2) {(b) 热路等价};
      
      % Heat source (power dissipation)
      \draw[therm_style, fill=orange!30] (0.5, 0.5) rectangle (1.5, 1.5);
      \node[align=center] at (1, 1) {热源\\$P_{\mathrm{heat}}$};
      
      % Thermal resistance (epitaxial layers)
      \draw[therm_style] (1.5, 1.0) -- (2.5, 1.0);
      \draw[therm_style, decorate, decoration={zigzag, segment length=4mm, amplitude=2mm}] 
        (2.5, 1.0) -- (4.5, 1.0) node[midway, above] {$R_{\mathrm{th}}$};
      \draw[therm_style] (4.5, 1.0) -- (5.5, 1.0);
      
      % Heat sink (boundary)
      \draw[therm_style, fill=blue!20] (5.5, 0.5) rectangle (6.5, 1.5);
      \node[align=center] at (6, 1) {热沉\\$T_{\mathrm{ambient}}$};
      
      % Temperature difference arrow
      \draw[<->, thick, red] (1.0, 1.7) -- (6.0, 1.7) node[midway, above] {$\Delta T = T_{\mathrm{junction}} - T_{\mathrm{ambient}}$};
      
      % Fourier's law box
      \draw[rounded corners, fill=red!15] (7.2, 2.5) rectangle (10.5, 3.8);
      \node[anchor=west, font=\small] at (7.4, 3.4) {傅里叶定律：};
      \node[anchor=west, font=\normalsize] at (7.4, 2.9) {$\Delta T = P \cdot R_{\mathrm{th}}$};
      
      % Labels
      \node[anchor=east, font=\small] at (0.3, 1.0) {功率耗散};
      \node[anchor=west, font=\small] at (6.7, 1.0) {恒温边界};
    \end{scope}
    
    % ========== Middle: Correspondence Table ==========
    \begin{scope}[shift={(0,-2)}]
      \draw[rounded corners, fill=gray!10] (-0.5,-2.8) rectangle (18.5, 0.3);
      
      % Table header
      \node[anchor=west, font=\small\bfseries] at (-0.3, 0.0) {电 - 热类比对照表：};
      
      % Correspondence rows
      \node[anchor=west, font=\small] at (-0.3, -0.5) {
        \begin{tabular}{ll@{\hspace{2em}}ll}
          \textbf{电学量} & \textbf{热学量} & \textbf{单位} \\
          \hline
          电压$V$ & 温差$\Delta T$ & V vs. K \\
          电流$I$ & 热流$P$ & A vs. W \\
          电阻$R$ & 热阻$R_{\mathrm{th}}$ & $\Omega$ vs. K/W \\
          电导率$\sigma$ & 热导率$\kappa$ & S/m vs. W/(m·K) \\
          电容$C$ & 热容$C_{\mathrm{th}}$ & F vs. J/K \\
        \end{tabular}
      };
      
      % Key formulas
      \node[anchor=west, font=\small] at (9, -0.5) {
        \textbf{热阻计算公式（一维近似）:}
      };
      \node[anchor=west, font=\normalsize] at (9, -1.2) {
        $\displaystyle R_{\mathrm{th}} = \frac{L}{\kappa A}$
      };
      \node[anchor=west, font=\footnotesize] at (9, -1.8) {
        其中$L$为传热路径长度，$\kappa$为材料热导率，$A$为截面积
      };
      
      % Typical values
      \node[anchor=west, font=\small] at (9, -2.4) {
        \textbf{典型值参考:} GaAs $\kappa \approx 46\,\mathrm{W/(m\cdot K)}$, InP $\kappa \approx 68\,\mathrm{W/(m\cdot K)}$
      };
    \end{scope}
    
    % ========== Application to PCSEL ==========
    \begin{scope}[shift={(0,-5)}]
      \draw[rounded corners, fill=yellow!15] (-0.5,-1.8) rectangle (18.5, 0.3);
      \node[anchor=west, font=\small] at (-0.3, -0.2) {
        \textbf{PCSEL中的应用示例}: 假设有源区功耗$P_{\mathrm{heat}} = 0.5\,\mathrm{W}$，外延层总热阻$R_{\mathrm{th}} = 80\,\mathrm{K/W}$，则温升
      };
      \node[anchor=west, font=\normalsize] at (-0.3, -0.9) {
        $\displaystyle \Delta T = P_{\mathrm{heat}} \cdot R_{\mathrm{th}} = 0.5 \times 80 = 40\,\mathrm{K}$
      };
      \node[anchor=west, font=\small] at (-0.3, -1.5) {
        这意味着结温$junction temperature$比热沉温度高$40^\circ\mathrm{C}$。若环境温度$25^\circ\mathrm{C}$，则结温达$65^\circ\mathrm{C}$，需要考虑温度对增益和折射率的影响。
      };
    \end{scope}
    
\end{tikzpicture}%
}
